DNS Configuration and Testing (20 Marks)

Step One: Verification of DNS Role Installation:

\begin{enumerate}
\def\labelenumi{\arabic{enumi}.}
\item
  Open Server Manager on the server.
\item
  Click the tools icon at the top right hand corner and select DNS.
\item
  Verify that your server appears in the console tree (GROUP1-SVR).
\end{enumerate}

\includegraphics[width=6.01088in,height=3.79942in]{media/image1.png}

Step Two: Configuring the Forward Lookup Zone

2.1. In the DNS Manager expand the GROUP1-SVR node.

2.2. Expand the Forward Lookup Zones.

2.3. Verify that the primary zone has been created for your domain i.e.
Group1.local this will have occurred during the AD DS configuration.

2.4.Click the zone to view the automatically created records (SOA, NS \&
A records).

\includegraphics[width=6in,height=4.13889in]{media/image2.png}

Step Three: Reverse Lookup Zone.

3.1. Right click on the Reverse Lookup Zones and select New Zone.

3.2. Click Next.

3.3. Select Primary zone and make sure Store the zone in Active
Directory is ticked. Hit next.

3.4. Select All DNS servers running on domain controllers in this
domain.

3.5. Select IPv4 Reverse Lookup Zone and hit next.

3.6. Enter your Network ID (i.e. 192.168.64.2) which will be the static
ip address on your Windows Server. Hit next.

3.7. Make sure Only Secure Dynamic Updates are allowed (Recommended for
Active Directory. Hit Finish.

\includegraphics[width=6in,height=3.63889in]{media/image3.png}

Step Four: Adding and Verifying DNS Records.

4.1. Navigate back to your Forward Lookup Zone (Group1cput.local).

4.2. Find your Host (A) record for Group {[}X{]} Windows Server 2019.
Double click to open properties.

4.3. Tick the box stating `Update associated pointer (PTR) record'. Hit
Apply and OK.

4.4. Go to your new Reverse Lookup Zone and refresh. Verify that there
is now a PTR record pointing to your server.

\includegraphics[width=6in,height=3.94444in]{media/image4.png}

Step Five: Testing DNS Resolution (Server \& Client)

Server Side-Test

5.1. Open PowerShell (Or Command Prompt if there are no restrictions on
admins).

5.2. Type `nslookup' and press enter. The Default Server should display
your server name and IP address.

5.3. Perform a forward lookup by entering your servers Fully Qualified
Domain Name (FQDN) (E.g. Group1WindowsServer2019.Group1cput.local) and
pressing enter. You should receive your server IP address.

5.4. Perform a reverse lookup by entering your server IP address (E.g.
192.168.64.2) and pressing enter. You should receive the FQDN.

\includegraphics[width=6in,height=4.68194in]{media/image5.png}

Client Side-Test

5.5. Switch to your client virtual machine running on your Semester 1
Hypervisor. Make sure it is on the same virtual network.

5.6. Open up command prompt/powerShell. Type ping Group1cput.local or
ping your server by name (ping Group1WindowsServer2019).

5.7. Verify that your client can resolve the name to the server IP
address and receive replies.

\includegraphics[width=6in,height=4.76944in]{media/image6.png}

References

McDonald, G., Papadopoulos, P., Pitropakis, N., Ahmad, J., \& Buchanan,
W. J. (2022). Ransomware: Analysing the impact on Windows Active
Directory Domain Services. Sensors, 22, 953. 69

Mokhtar, B., Jurcut, A., ElSayed, M., \& Azer, M. (2022). Active
Directory attacks---Steps, types, and signatures.Electronics , 11, 2629.
41

Moric, Z., Dakic, V., Redzepagic, J., \& Otocan, F. (2024). Penetration
testing and exploitation of Active Directory configuration
vulnerabilities. 2024 4th Intelligent Cybersecurity Conference (ICSC) ,
118-126. 4

Pongsrisomchai, S., \& Ngamsuriyaroj, S. (2019). Automated IT audit of
Windows Server access control. 2019 21st International Conference on
Advanced Communication Technology (ICACT) , 539-544. 1
